\section*{Control de cambios respecto a la entrega antigua}
\phantomsection
\addcontentsline{toc}{section}{Control de cambios respecto a la entrega antigua}
\begingroup
\small
\noindent\textbf{Comparación:} \texttt{PROPUESTA\_TECNICA\_ANTIGUO.pdf} frente a los fuentes de \texttt{DGP\_P2\_buena}.

Cada código PT tiene un color propio, aplicado también al contenido nuevo o modificado que se indica. Las figuras se identifican con un marco del mismo color, conservando la imagen original. El texto que no cambia mantiene su formato habitual. Las supresiones se describen en el registro al no existir un bloque actual que colorear.

Este registro describe diferencias efectivas entre versiones. La tabla preexistente de observaciones CC-01 a CC-08 se conserva como contenido de la versión nueva; sus propuestas pendientes no se consideran cambios ya realizados.

\entradaCambio{PT-01}{cambioPresentacion}{Presentación del documento}
{La portada y el encabezado utilizaban la presentación anterior; varias secciones principales compartían página.}
{Se aplica la paleta CLENCH, se incorpora un logotipo central a la portada y un logotipo en la cabecera, y se introducen saltos de página antes de las secciones principales.}
{Portada, cabeceras y comienzo de las secciones; títulos de portada y marco del logotipo central marcados en gris pizarra.}

\entradaCambio{PT-02}{cambioRevision}{Observaciones de revisión incorporadas}
{El documento pasaba directamente del índice al resumen, sin la tabla de observaciones de revisión.}
{Se añade el bloque «Control de Cambios», con ocho observaciones CC-01 a CC-08 sobre enfoque al cliente, requisitos, solución, beneficios, cronograma, documentación, seguimiento y maquetación, indicando su estado y alcance.}
{Bloque preexistente anterior al resumen, identificado en su totalidad con el color de PT-02.}

\entradaCambio{PT-03}{cambioObjetivos}{Objetivos orientados al alumnado y trazabilidad}
{El objetivo general mencionaba una interfaz accesible. Los específicos trataban alcance, personalización, implementación del MVP, verificación y preparación de la entrega, sin enlaces a requisitos.}
{El objetivo general añade la personalización. Se reformulan los cinco objetivos para tratar alcance, interacción personalizable, aprendizaje autónomo, adaptación individual y fiabilidad; se añaden relaciones y enlaces a RF/RNF y precisiones sobre validación con el centro, requisitos pospuestos y disponibilidad no comprometida.}
{Sección 2 completa. El resumen, el tipo de caso de negocio y la duración se mantienen.}

\entradaCambio{PT-04}{cambioCatalogo}{Presentación del catálogo funcional}
{La cabecera decía «Requisito propuesto» y RF-17 tenía la prioridad «MVP (revisar)».}
{La cabecera pasa a «Requisito funcional» y la prioridad de RF-17 queda como «MVP». Se mantienen los 41 requisitos funcionales y sus criterios de aceptación.}
{Cabeceras repetidas de la tabla de 3.1 y celda de prioridad de RF-17.}

\entradaCambio{PT-05}{cambioRNF}{Requisitos no funcionales verificables}
{Los ocho requisitos no funcionales se presentaban en una lista temática, sin identificadores RNF ni columna propia de criterios de aceptación.}
{Se presentan en una tabla como RNF-01 a RNF-08, con el requisito y los pasos de aceptación para accesibilidad, idioma, privacidad, uso sin conexión, usabilidad, adaptación de interfaz, consistencia y rendimiento.}
{Apartado 3.2 completo.}

\entradaCambio{PT-06}{cambioInformacion}{Catálogo de requisitos de información}
{No existía un apartado específico de requisitos de información; los elementos pospuestos ocupaban el apartado 3.3.}
{Se añade una tabla INF-01 a INF-08 con cuentas, estudiante y accesibilidad, avatares, asociaciones, tareas, alternativas, tipos de tareas e instrucciones, y sus atributos. Los elementos pospuestos pasan a 3.4 sin cambiar su lista.}
{Nuevo apartado 3.3 completo.}

\entradaCambio{PT-07}{cambioArquitectura}{Explicación y representación de la solución}
{Se describía la aplicación, un esquema lógico tabular, un listado de componentes y un flujo principal en cinco pasos.}
{Se relaciona la propuesta provisional con la visita al centro y con la validación pendiente; el esquema tabular se sustituye por una imagen. El flujo pasa a cuatro pasos con enlaces a RF/RNF: acceso, planificación, consulta accesible y realización/continuidad. El listado independiente de componentes deja de aparecer en esta propuesta.}
{Sección 6 completa, incluida la figura. En el manual nuevo, el apartado 3.2 incorpora la descripción de los componentes.}

\entradaCambio{PT-08}{cambioBeneficios}{Beneficios y perfiles beneficiarios}
{La tabla recogía cuatro grupos: alumnado, profesionales y administración agrupados, familias y centro.}
{Se amplía a nueve perfiles, separando profesorado, apoyo educativo, orientación y administración e incorporando futuras promociones y comunidad universitaria. Se detallan alcance, beneficios, verificación y trazabilidad; se añaden condiciones de validación, mantenimiento y privacidad, y posibles indicadores y una línea base.}
{Apartado 9.2 completo. La tabla de valor añadido de 9.1 se mantiene.}

\entradaCambio{PT-09}{cambioDocumentacion}{Documentación y actas}
{Tras los registros internos se explicaban las pautas de actas y su nomenclatura, y se incluía el apartado 10.3 con una plantilla de acta completa.}
{El cierre se resume en un párrafo sobre repositorio, frecuencia de actualización y revisión. Desaparecen de la propuesta las pautas detalladas y la plantilla. Se mantienen las tablas DC-01 a DC-06 y RI-01 a RI-06; el manual nuevo incorpora pautas de actas y actas cumplimentadas.}
{Párrafo final de 10.2; las eliminaciones quedan documentadas en esta entrada.}

\entradaCambio{PT-10}{cambioCronograma}{Anexo visual del cronograma}
{El anexo A contenía una sola figura de planificación de historias de usuario por sprint.}
{Se sustituye por cuatro imágenes del cronograma, con cuatro pies de figura que figuran en los fuentes como «Diagrama de Arquitectura (Parte 1--4)». El texto de Actividades y Cronograma se mantiene.}
{Anexo A: título, marcos y pies de las cuatro figuras; se conserva la denominación original de sus pies.}
\endgroup
